%--------------------------------------------------------------------------------
%
%
%--------------------------------------------------------------------------------
\pagebreak
\section*{NOP — No Operation}
\addcontentsline{toc}{section}{NOP — No Operation}

\begin{iPageDescEntry}{Syntax}
    \texttt{NOP}
\end{iPageDescEntry}

\begin{iPageDescEntry}{Format}
    The \texttt{NOP} instruction uses no format. \\

    \begin{flushleft}
    \drawInstrFormat
        {0,1,3,16}
        {\OpCodeGrpALU, \OpCodeNOP, 0}
    \end{flushleft}
\end{iPageDescEntry}

\begin{iPageDescEntry}{Description}
    The \texttt{NOP} (No Operation) instruction does nothing. 
    It is typically used for timing, alignment, or to occupy a pipeline slot.
\end{iPageDescEntry}

\begin{iPageDescEntry}{Operation}
    \texttt{// no changes to registers or memory}
\end{iPageDescEntry}

\begin{iPageDescEntry}{Exceptions}
    None.
\end{iPageDescEntry}

\begin{iPageDescEntry}{Notes}
    - Often used for delay loops or to align instructions in memory. \\
    - Safe to execute in any mode, does not trigger traps or exceptions.
\end{iPageDescEntry}

